% BHU PYQ -- Manually transcribed & error-corrected
% Paper: MTB-AM-203 : Ancillary Mathematics I
% Exam: B.Sc. (Hons.) Semester II Examination, 2013-14

\documentclass[11pt,a4paper]{article}
\usepackage[a4paper,top=2.5cm,bottom=2.5cm,left=2.8cm,right=2.8cm,headheight=14pt]{geometry}
\usepackage{amsmath,amssymb}
\usepackage[T1]{fontenc}\usepackage[utf8]{inputenc}
\usepackage{lmodern,microtype,fancyhdr,enumitem,hyperref,lastpage,parskip}

\hypersetup{colorlinks=true,urlcolor=black,linkcolor=black,
  pdftitle={MTB-AM-203 Ancillary-I -- BHU B.Sc. Sem II 2013-14}}
\pagestyle{fancy}\fancyhf{}
\renewcommand{\headrulewidth}{0.4pt}\renewcommand{\footrulewidth}{0.4pt}
\fancyhead[L]{\small\textit{Banaras Hindu University}}
\fancyhead[R]{\small\textit{MTB-AM-203: Ancillary Mathematics I}}
\fancyfoot[C]{\small Page \thepage\ of \pageref{LastPage}}

\newlist{parts}{enumerate}{2}
\setlist[parts,1]{label=(\alph*),leftmargin=*,itemsep=5pt,topsep=3pt}
\setlist[parts,2]{label=(\roman*),leftmargin=*,itemsep=3pt,topsep=2pt}
\newcommand{\pts}[1]{\hfill\textit{[#1]}}

\begin{document}
\begin{center}
  {\large\bfseries BANARAS HINDU UNIVERSITY}\\[2pt]
  {\normalsize B.Sc.\ (Hons.) --- Semester II Examination, 2013-14 \textnormal{(Ancillary)}}\\[6pt]
  \rule{0.8\linewidth}{0.5pt}\\[6pt]
  {\large\bfseries Paper No.\ MTB-AM-203: Ancillary Mathematics I}\\[4pt]
  {\normalsize Time: 3 Hours \hfill Full Marks: 70}\\[2pt]
  {\small Code:\ BS/Sem\,II/4}
\end{center}
\rule{\linewidth}{0.4pt}
\smallskip
\noindent\textit{Instructions: Answer \textbf{five} questions including Question~1 which is compulsory. Figures in right-hand margin indicate marks.}
\bigskip

\noindent\textbf{Question 1.} \textit{(Compulsory --- 2 marks each.)}
\begin{parts}
  \item Which of the sets are equal: (i)~$\{x:x$ is a letter in `follow'$\}$; (ii)~$\{x:x$ is a letter in `wolf'$\}$; (iii)~$\{x:x$ is a letter in `flow'$\}$.
  \item If $A=\{3,4\}$ and $B=\{1,2,5\}$, find $A\times B$.
  \item If $A=\{1,2,3,4\}$, $B=\{2,4,6,8\}$, $C=\{3,4,5,6\}$, find $(A\cap B)\cap C$.
  \item Which of the following is a function $f:X\to X$ for $X=\{1,2,3,4\}$: $f_1=\{(2,3),(1,4),(2,1),(3,2),(4,4)\}$? (State reason.)
  \item $R=\{(x,y):x+2y=5\}$ on $\mathbb{N}$. Is $R$ transitive?
  \item Find $k$ if $\begin{vmatrix}5&k+2\\k-2&3\end{vmatrix}=\begin{vmatrix}k+3&4\\-k&2\end{vmatrix}$.
  \item Find $AB$ where $A=\bigl[\begin{smallmatrix}3&2&0\end{smallmatrix}\bigr]$, $B=\bigl[\begin{smallmatrix}0\\1\end{smallmatrix}\bigr]$.
  \item If $A=\bigl[\begin{smallmatrix}1&2\\4&1\end{smallmatrix}\bigr]$ and $B=\bigl[\begin{smallmatrix}2&1\\-1&0\end{smallmatrix}\bigr]$, find $A-2B$.
  \item Without expanding the determinant, prove $\begin{vmatrix}a&b&c\\x&y&z\\p&q&r\end{vmatrix}=\begin{vmatrix}x&y&z\\p&q&r\\a&b&c\end{vmatrix}$.
  \item Determine the rank of $\begin{pmatrix}1&2&3\\2&4&6\end{pmatrix}$.
  \item If $n(A\cup B)=20$, $n(A\cap B)=8$, $n(B)=12$, find $n(A)$.
\end{parts}
\medskip\rule{\linewidth}{0.2pt}

\medskip\noindent\textbf{Question 2.}
\begin{parts}
  \item Show $n(A\cup B\cup C)=n(A)+n(B)+n(C)-n(A\cap B)-n(A\cap C)-n(B\cap C)+n(A\cap B\cap C)$.\pts{6}
  \item In a survey: 1051 used bus, 709 car, 1252 train, 406 bus and car, 316 bus and train, 301 train and car, 51 all three. How many were surveyed?\pts{6}
\end{parts}
\medskip\rule{\linewidth}{0.2pt}

\medskip\noindent\textbf{Question 3.}
\begin{parts}
  \item Evaluate $\begin{vmatrix}167&19&21\\39&13&14\\81&24&26\end{vmatrix}$.\pts{6}
  \item Show $\begin{vmatrix}1&1&1\\a&b&c\\a^2&b^2&c^2\end{vmatrix}=(a-b)(b-c)(c-a)$.\pts{6}
\end{parts}
\medskip\rule{\linewidth}{0.2pt}

\medskip\noindent\textbf{Question 4.}
\begin{parts}
  \item Solve by determinants: $3x+5y-7z=13$, $4x+y-12z=6$, $2x+9y-3z=20$.\pts{6}
  \item (i) Find $x$ if $\begin{vmatrix}2x&0\\x&x\end{vmatrix}=\begin{vmatrix}10&0\\-1&2\end{vmatrix}$. (ii) Solve for $X$: $X\begin{pmatrix}2&5\\1&4\end{pmatrix}=\begin{pmatrix}-6\\0\end{pmatrix}$.\pts{3+3}
\end{parts}
\medskip\rule{\linewidth}{0.2pt}

\medskip\noindent\textbf{Question 5.}
\begin{parts}
  \item If $A=\begin{pmatrix}2&1&-1\\1&1&2\\7&2&1\end{pmatrix}$ and $B=\begin{pmatrix}0&2&1\\1&3&0\\-9&1&2\end{pmatrix}$, compute $A^2-4B^2$.\pts{6}
  \item If $A=\begin{pmatrix}1&1&2\\1&3&2\\1&-1&2\end{pmatrix}$, $B=\begin{pmatrix}0&2&1\\1&0&0\\5&0&1\end{pmatrix}$, $C=\begin{pmatrix}4\\4\\1\end{pmatrix}$, compute $(AB)C$ and $A(BC)$ and verify $(AB)C=A(BC)$.\pts{6}
\end{parts}
\medskip\rule{\linewidth}{0.2pt}

\medskip\noindent\textbf{Question 6.}
\begin{parts}
  \item Find the inverse of $A=\begin{pmatrix}0&1&-2\\-1&0&5\\2&-5&0\end{pmatrix}$.\pts{6}
  \item Find the rank of $A=\begin{pmatrix}0&2&2&-1\\1&0&4&1\\3&1&0&9\\1&1&-2&9\end{pmatrix}$.\pts{6}
\end{parts}
\medskip\rule{\linewidth}{0.2pt}

\medskip\noindent\textbf{Question 7.}
\begin{parts}
  \item Solve: $x+3y-2z=0$, $2x-y+4z=0$, $x-11y+14z=9$.\pts{6}
  \item Test consistency and solve: $x+2y+3z=14$, $3x+y+2z=11$, $2x+3y+z=11$.\pts{6}
\end{parts}

\vfill\rule{\linewidth}{0.4pt}
\begin{center}\small
\href{https://bhu-exam.vercel.app/}{Click here to visit our website}\\[4pt]
\href{https://me-aryan.vercel.app/}{Click here to visit our contributor}\\[4pt]
\href{https://quashan-me.vercel.app/}{Click here to visit our contributor}
\end{center}
\end{document}
